\chapter{What PebbleIR Must Accomplish}
\label{ch:pebble-purpose}

A parser can build a beautiful tree that no optimizer wants to edit.
An optimizer can build a compact graph that no interpreter can execute unambiguously.
An \term{intermediate representation} succeeds when these clients share a precise account of what the graph means and how it may change.
PebbleIR is an exercise in making that account executable.

\section{Start with a client}
\label{sec:pebble-client}
Our first client translates \code{(x+7)*(3+5)}.
It needs literals, an input value, addition, multiplication, and a return.
A second client folds the literal sum.
A third executes the result for every permitted input.
Their first disagreement should be diagnosed as a contract violation rather than hidden by an undocumented convention.

The value domain is the integers from zero through 255, with addition and multiplication reduced modulo 256.
At input 9 the expression returns 128.
At input 250, the intermediate sum becomes 1 and the final result is 8.
The wraparound case belongs in the first test, because a host-language implementation that accidentally uses different arithmetic can otherwise appear correct.

\section{A deliberately small agreement}
\label{sec:pebble-scope}
Table~\ref{tab:pebble-scope} states the implementation boundary.
A limitation belongs in the interface when a caller must plan around it.
In particular, exceeding capacity is a reported construction failure, not permission to overwrite an array.

\begin{table}[htbp]
\caption{The first implementation chooses explicit limits and semantics.}
\label{tab:pebble-scope}
\centering
\begin{tabular}{ll}
\toprule
Property & Contract \\
\midrule
Function & One byte input and one byte return \\
Values & Byte integers and booleans \\
Control & Return, jump, conditional branch \\
Merges & Two-input phis indexed by predecessor block \\
Storage & Eight blocks, 128 value slots \\
Effects & No loads, stores, calls, or exceptions \\
Execution & Loops allowed; finite fuel is a harness policy \\
\bottomrule
\end{tabular}
\end{table}

The absence of effects makes dead arithmetic removable, but does not make a diverging loop removable.
The absence of pointers removes alias analysis from this implementation, but does not make alias analysis unnecessary in a future memory extension.
The interface should expose these boundaries without pretending that a small representation is a universal compiler architecture.

\section{Establish the shortest complete path}
\label{sec:pebble-first-path}
Algorithm~\ref{alg:pebble-first-path} is the acceptance procedure for the first milestone.
It is implemented by the expression compiler, interpreter, optimizer, and VM tests introduced in later chapters.
The source oracle computes from a separate expression tree; it does not ask the optimizer what the answer should be.

\begin{algorithm}[htbp]
\caption{Accepting an end-to-end PebbleIR specimen.}
\label{alg:pebble-first-path}
\begin{pseudocode}
\Require expression tree $e$ over byte arithmetic, sufficient construction capacity
\Ensure checked specimen or a reproducible disagreement
\State compile its printed expression into $F$; verify $F$
\State copy $F$ into $G$ and optimize $G$; verify $G$
\State lower $G$ to virtual-machine code $M$
\ForAll{$x\in\{0,\ldots,255\}$}
\State compare tree evaluation, $F$, $G$, and $M$ under sufficient fuel
\State report the first input whose results disagree
\EndFor
\end{pseudocode}
\end{algorithm}

For a tree of $N$ nodes this performs 256 bounded executions, not an exhaustive proof over all possible trees.
That distinction directs the next investment: prove the local arithmetic rules, then generate structurally different programs to test their implementation.
Do not add five optimization passes before one complete translation path works.
An incomplete chain gives every component an excuse to blame a component that does not exist yet.

\section{Implementation checkpoint}
\label{sec:pebble-checkpoint}
The repository places the implementation under \code{example-code/pebbleir/} so the pamphlet and its code travel together.
The header is the public model; the C file supplies its operations; the test file is a client; the demo is a second client with a smaller interface.
The existing third-party directories are not dependencies of this C11 specimen.
An allocator or container library can be adopted later when a measured requirement justifies the extra interface.

\section{Exercises}
\label{sec:pebble-one-exercises}
\begin{enumerate}
\item \label{ex:pebble-wrap} Evaluate the running expression at inputs 0, 9, and 250, showing every modular reduction.
\item \label{ex:pebble-capacity-contract} A builder has one unused value slot and is asked to build a binary expression with two new constants. Explain why checking only the final operation's allocation is insufficient, and propose a failure contract.
\end{enumerate}

\section{Further reading}
\label{sec:pebble-one-reading}
Appel places intermediate representations inside a complete translation path; read the neighboring phases as well as the representation itself~\cite{appel-1998-modern}.
Safonov provides a complementary perspective on explicit compiler trust obligations~\cite{safonov-2010-trustworthy}.
